\documentclass[11pt]{article}
\usepackage[a4paper,margin=1in]{geometry}
\usepackage{amsmath,amssymb,amsthm,mathtools}
\usepackage{booktabs}
\usepackage{graphicx}
\usepackage{microtype}
\usepackage[hidelinks]{hyperref}
\usepackage{xcolor}

\newtheorem{theorem}{Theorem}
\newtheorem{lemma}{Lemma}
\newcommand{\R}{\mathbb{R}}
\newcommand{\F}{\mathcal{F}}
\newcommand{\SN}{\mathcal{S}_N}
\newcommand{\dd}{\,\mathrm{d}}

\title{An Exact Dimension Threshold for the Localized\\
Aubin--Talenti Family in arXiv:2504.15968}
\author{Candidate full solution produced in run \texttt{fa\_banach\_001}\\
Model: GPT5.6}
\date{19 July 2026}

\begin{document}
\maketitle

\begin{center}
\fbox{\parbox{0.91\textwidth}{\textbf{Status.}
Candidate full solution, currently assessed as likely valid and awaiting
human review. The novelty search described below was bounded, not exhaustive.}}
\end{center}

\section{Source and open question}

The source is S. Chakraborty, D. Gupta, S. Malhotra, and K. Sreenadh,
\emph{Global Compactness Result for a Br\'ezis--Nirenberg-Type Problem
Involving Mixed Local Nonlocal Operator}, arXiv:2504.15968v2 (11 May
2025)~\cite{source}. On page 7, Remark 1.1(1), the authors ask for the
threshold dimension beyond which the energy of their localized
Aubin--Talenti family (1.7) is strictly below $2^{2/N}\SN$.

\begin{figure}[ht]
\centering
\includegraphics[width=\textwidth]{figures/open_problem_crop.png}
\caption{The open question in arXiv:2504.15968v2, page 7,
Remark 1.1(1). The crop retains the full readable page width.}
\end{figure}

In the paper's notation, the unnormalized Gagliardo seminorm is
\[
 [u]_s^2=\iint_{\R^{2N}}
 \frac{|u(x)-u(y)|^2}{|x-y|^{N+2s}}\dd x\dd y,
 \qquad 0<s<1,
\]
and the mixed Rayleigh quotient is
\[
 S(u;\R^N)=
 \frac{\|\nabla u\|_2^2+[u]_s^2}{\|u\|_{2^*}^2},
 \qquad 2^*=\frac{2N}{N-2}.
\]
The standard bubble is
\[
 U_N(x)=(1+|x|^2)^{-(N-2)/2}.
\]

\section{Proof intuition}

The source estimates $[U_N]_s$ by interpolating between the $L^2$ and
gradient norms. That is sufficient asymptotically but loses the dimension at
which the inequality first becomes true. Here the seminorm is evaluated
exactly. The unitary Fourier transform of $U_N$ is a multiple of
$|\xi|^{-1}K_1(|\xi|)$, so the fractional energy is a Mellin integral of
$K_1^2$ and hence a product of Gamma functions. After division by the exact
gradient energy, the source question becomes one explicit scalar inequality.

The remaining issue is whether success can later revert to failure as $N$
grows. The quotient between dimensions $N+2$ and $N$ collapses to an
elementary rational expression. Once normalized by the critical barrier
$2^{2/N}-1$, this quotient is strictly less than one for $N\geq7$. Thus the
two parity classes are monotone: two consecutive successful dimensions certify
every larger dimension and give a finite exact stopping rule.

\section{Exact theorem}

For $N\geq5$ and $0<s<1$, set
\begin{align}
 Q_{N,s}
 &:=\frac{[U_N]_s^2}{\|\nabla U_N\|_2^2} \notag\\
 &=\frac{
 2\pi^{N/2}|\Gamma(-s)|
 \Gamma(\frac N2+s-1)^2
 \Gamma(\frac N2+s-2)\Gamma(N)}{
 N(N-2)\Gamma(\frac N2-1)^2\Gamma(\frac N2)^2
 \Gamma(N+2s-2)},                                      \label{eq:Q}
\end{align}
and
\[
 h_N:=2^{2/N}-1,\qquad F_{N,s}:=\frac{Q_{N,s}}{h_N}.
\]

Let $\mathsf P_N(s)$ denote the precise uniform large-cutoff property used in
the source's Lemma 4.3: there are $S_1<2^{2/N}\SN$ and $R_0$ such that the
normalized localized family (1.7) has Rayleigh quotient at most $S_1$ whenever
$R\geq R_0$, uniformly in $t\in[0,1)$ and $\sigma\in\mathbb S^{N-1}$.

\begin{theorem}[Exact threshold and finite stopping rule]\label{thm:main}
For each $s\in(0,1)$ the following statements hold.
\begin{enumerate}
\item For every $N\geq5$,
\[
 \mathsf P_N(s)\quad\Longleftrightarrow\quad
 Q_{N,s}<2^{2/N}-1
 \quad\Longleftrightarrow\quad F_{N,s}<1.
\]
\item For every $N\geq7$,
\[
 F_{N+2,s}<F_{N,s}.
\]
Consequently both parity subsequences are strictly decreasing from dimension
7 onward and tend to zero.
\item The exact eventual threshold is
\[
 N_{\mathrm{th}}(s)
 =1+\max\{N\geq5:F_{N,s}\geq1\}.                       \label{eq:nth}
\]
The set in \eqref{eq:nth} is finite. It is determined by a finite procedure:
evaluate the explicit formula \eqref{eq:Q} in increasing dimension, retaining
the earlier failures, and stop after two consecutive successes beginning at
a dimension at least 7. No later failure is possible, so one plus the last
retained failure is exactly $N_{\mathrm{th}}(s)$.
\end{enumerate}
\end{theorem}

Thus the question is reduced without estimates to a closed Gamma-function
test, together with a proved finite stopping rule. For illustration, the
included 80-digit checker gives the following values.
\begin{center}
\begin{tabular}{cccccccc}
\toprule
$s$ & $0.05$ & $0.10$ & $0.25$ & $0.50$ & $0.75$ & $0.90$ & $0.99$\\
\midrule
$N_{\mathrm{th}}(s)$ & 22 & 21 & 19 & 19 & 20 & 22 & 25\\
\bottomrule
\end{tabular}
\end{center}
These decimal evaluations are checks of the exact criterion, not ingredients
of the proof.

\section{Proof of Theorem \ref{thm:main}}

\subsection{Exact fractional energy}

Use the unitary Fourier transform
\[
 \widehat f(\xi)=(2\pi)^{-N/2}
 \int_{\R^N}e^{-ix\cdot\xi}f(x)\dd x.
\]
Write $B=N/2$ and $A=B+s-1$. The Gamma integral representation of
$(1+|x|^2)^{-(B-1)}$, followed by the Gaussian Fourier transform, gives
\begin{equation}
 \widehat U_N(\xi)=
 \frac{2^{2-B}}{\Gamma(B-1)}\,
 \frac{K_1(|\xi|)}{|\xi|}.                             \label{eq:FU}
\end{equation}
Indeed the remaining Laplace integral is
\[
 \int_0^\infty t^{-2}e^{-t-|\xi|^2/(4t)}\dd t
 =\frac{4K_1(|\xi|)}{|\xi|}.
\]

Let
\[
 c_{N,s}=\frac{4^s\Gamma(B+s)}{
 \pi^B|\Gamma(-s)|}.
\]
This is the multiplier normalization for the fractional Laplacian, and
Plancherel gives the raw-seminorm identity
\begin{equation}
 [f]_s^2=\frac{2}{c_{N,s}}
 \int_{\R^N}|\xi|^{2s}|\widehat f(\xi)|^2\dd\xi.
                                                               \label{eq:raw}
\end{equation}
The standard Mellin integral~\cite[Section 10.43]{dlmf}, valid for
$\Re\mu>2|\Re\nu|$, is
\begin{equation}
 \int_0^\infty r^{\mu-1}K_\nu(r)^2\dd r
 =\frac{2^{\mu-3}}{\Gamma(\mu)}
 \Gamma\!\left(\frac{\mu}{2}\right)^2
 \Gamma\!\left(\frac{\mu}{2}+\nu\right)
 \Gamma\!\left(\frac{\mu}{2}-\nu\right).            \label{eq:Ksquare}
\end{equation}
Apply \eqref{eq:Ksquare} with $\mu=N+2s-2=2A$ and $\nu=1$.
Equations \eqref{eq:FU}--\eqref{eq:raw}, together with
$\omega_{N-1}=2\pi^B/\Gamma(B)$, yield
\begin{equation}
 [U_N]_s^2=
 \frac{\omega_{N-1}\pi^B|\Gamma(-s)|
 \Gamma(A)^2\Gamma(A-1)}{
 \Gamma(B-1)^2\Gamma(2A)}.                              \label{eq:rawU}
\end{equation}
All conditions are satisfied for $N\geq5$ and $0<s<1$.

A radial beta integral gives
\begin{equation}
 \|\nabla U_N\|_2^2
 =\frac{\omega_{N-1}N(N-2)}{2}
 \frac{\Gamma(B)^2}{\Gamma(N)}.                        \label{eq:gradU}
\end{equation}
Dividing \eqref{eq:rawU} by \eqref{eq:gradU} proves
\eqref{eq:Q}.

\subsection{Equivalence with the localized family}

The source proves uniform convergence of its cut-off bubbles to the uncut
translated/dilated family in the gradient norm, the Gagliardo seminorm, and
the critical norm~\cite[Section 4]{source}. Translation and critical scaling
give
\begin{align*}
 \|\nabla U[t\sigma,1-t]\|_2^2&=\|\nabla U_N\|_2^2,\\
 \|U[t\sigma,1-t]\|_{2^*}^2&=\|U_N\|_{2^*}^2,\\
 [U[t\sigma,1-t]]_s^2&=(1-t)^{2-2s}[U_N]_s^2.
\end{align*}
Since $U_N$ attains the sharp classical Sobolev quotient,
\[
 S(U[t\sigma,1-t];\R^N)
 =\SN\bigl(1+(1-t)^{2-2s}Q_{N,s}\bigr).
\]
The supremum over $t\in[0,1)$ occurs at $t=0$. Therefore its limiting value
is below $2^{2/N}\SN$ exactly when $Q_{N,s}<h_N$. A strict inequality leaves a
positive gap, and uniform cutoff convergence supplies the required $S_1$ and
$R_0$. Conversely, if $Q_{N,s}\geq h_N$, the $t=0$ cutoffs converge to a value
at or above the threshold, so no fixed $S_1<2^{2/N}\SN$ can work for all large
$R$. This proves part 1.

\subsection{Monotonicity in dimension}

Gamma recurrence in \eqref{eq:Q} simplifies completely:
\begin{equation}
 \frac{Q_{N+2,s}}{Q_{N,s}}
 =\frac{2\pi(N+1)(N+2s-2)(N+2s-4)}
 {(N+2s-1)N(N-2)(N+2)}.                                \label{eq:rec}
\end{equation}
For $N\geq7$, the factor
$x(x-2)/(x+1)$ is increasing at $x=N+2s-2$, hence the right side of
\eqref{eq:rec} is strictly below its $s\uparrow1$ limit:
\begin{equation}
 \frac{Q_{N+2,s}}{Q_{N,s}}<\frac{2\pi}{N+2}.           \label{eq:qbound}
\end{equation}

It remains to compare the barriers. Put $h_N=2^{2/N}-1$. For $N\geq7$,
$h_N\leq h_7<11/50$. The latter inequality follows by raising
$2^{2/7}<61/50$ to the seventh power, since $(61/50)^7>4$.
The classical rational bounds $\log2>173/250$ and $\pi<22/7$ imply
$\log2/\pi>11/50$. Consequently
\[
 2\pi h_N<2\log2<(N+2)h_{N+2},
\]
where the second inequality uses $e^x-1>x$ at
$x=2\log2/(N+2)$. Combining this with \eqref{eq:qbound} gives
\[
 \frac{F_{N+2,s}}{F_{N,s}}
 =\frac{Q_{N+2,s}}{Q_{N,s}}\frac{h_N}{h_{N+2}}<1.
\]
For completeness, $e^x-1<xe^x$ and $e^y-1>y$ also give
\[
 \frac{h_N}{h_{N+2}}<\frac{N+2}{N}\,2^{2/N}<2
 \qquad(N\geq7).
\]
Together with \eqref{eq:qbound}, this makes the quotient
$F_{N+2,s}/F_{N,s}$ of order $O(1/N)$, so both parity subsequences tend to zero.
This proves part 2, finiteness in \eqref{eq:nth}, and the stopping rule in
part 3.

\section{Verification notes}

The included script \texttt{code/threshold\_verifier.py} evaluates
\eqref{eq:Q} in log-Gamma form, verifies \eqref{eq:rec} against independent
Gamma evaluations, and checks representative thresholds. With 80 decimal
digits, the maximum relative recurrence discrepancy over the sampled cases
was below $2\times10^{-78}$. Direct numerical quadrature of the Bessel
integral at $(N,s)=(5,1/2),(8,1/4),(19,3/4)$ also agreed with
\eqref{eq:Q}.

There is a useful normalization check independent of those calculations.
Taking $s\uparrow1$ in \eqref{eq:rawU} gives
\[
 \lim_{s\uparrow1}(1-s)\frac{[U_N]_s^2}{\|\nabla U_N\|_2^2}
 =\frac{\omega_{N-1}}{2N},
\]
which is exactly the Bourgain--Brezis--Mironescu constant for the raw double
integral used by the source.

\section{Novelty search, scope, and limitations}

A bounded search was performed on 19 July 2026. It covered the exact arXiv id
2504.15968, the quoted phrase ``precisely determine the threshold dimension'',
``localized Aubin--Talenti bubbles'', ``mixed local nonlocal'', and the level
$2^{2/N}\SN$, using arXiv-focused and general web search. The current source
remains arXiv v2, last revised 11 May 2025~\cite{source}. The search found
related later work on mixed local--nonlocal critical problems but no paper
claiming this exact threshold calculation. Thus no known duplicate was found,
but the search is not an exhaustive bibliographic guarantee.

This result determines the threshold for the specific localized family and
the uniform fixed-gap statement used in the source. It does not claim that a
different test family cannot establish the PDE existence theorem in lower
dimensions, nor does it relax the domain assumptions in Remark 1.1(3).

\section{Human-review recommendation}

The result is suitable for expert review as a candidate full answer to Remark
1.1(1). The highest-value checks are: (i) the factor $2/c_{N,s}$ in
\eqref{eq:raw}; (ii) the cancellation leading to \eqref{eq:rawU}; and (iii)
that the source's uniform cutoff lemma is interpreted with a fixed
$S_1<2^{2/N}\SN$, so equality at the uncut limit is correctly classified as
failure.

\begin{thebibliography}{9}
\bibitem{source}
S. Chakraborty, D. Gupta, S. Malhotra, and K. Sreenadh,
\emph{Global Compactness Result for a Br\'ezis--Nirenberg-Type Problem
Involving Mixed Local Nonlocal Operator}, arXiv:2504.15968v2, 2025.

\bibitem{dlmf}
NIST Digital Library of Mathematical Functions,
\emph{Modified Bessel Functions: Integrals}, Section 10.43,
\url{https://dlmf.nist.gov/10.43}, accessed 19 July 2026.
\end{thebibliography}

\end{document}
